\chapter{Glossary}
\label{app:glosarium}

The terms below are used with fixed technical meanings throughout the book. Indonesian
equivalents are given in parentheses, because the Indonesian edition of this book uses those
terms and much of the local literature follows it.

\begin{description}[style=nextline,itemsep=0.35em]
  \item[Absolute convergence] a series that still converges when all its terms are forced positive (\emph{konvergensi mutlak}).
  \item[Absolute value] the distance of a number from zero on the number line, regardless of its sign (\emph{nilai mutlak}).
  \item[Accumulation point] a point every neighborhood of which contains a member of the set other than itself (\emph{titik akumulasi}).
  \item[Activation function] a nonlinear mapping on the output of a unit; it decides how much signal is passed on and the shape of its derivative (\emph{fungsi aktivasi}).
  \item[Adaptive attack] an attack designed with knowledge of the defense being tested, including knowledge that its gradient is manipulated (\emph{serangan adaptif}).
  \item[Adjoint] in the backward pass, the number that says how strongly the final loss is affected by one quantity in the middle of the network; one adjoint per quantity (\emph{adjoint}).
  \item[Adjugate] the cofactor matrix arranged in transpose; used to write the inverse as a single formula, $A^{-1}=\operatorname{adj}(A)/\det A$ (\emph{adjoin}).
  \item[Adversarial example] a small disturbance to the input, deliberately made to change the output of a model (\emph{serangan adversarial}).
  \item[Adversarial learning] the study of examples deliberately designed to fool a model, and of how to defend against them (\emph{pembelajaran adversarial}).
  \item[Antiderivative] a function whose derivative is the function being integrated; the way home from a derivative (\emph{antiturunan}).
  \item[Asymptotically stable] stays close after being disturbed, and then actually returns to the equilibrium point (\emph{stabil asimtotik}).
  \item[Axiom] a statement accepted without proof as a starting point; the whole content of a chapter rests on it (\emph{aksioma}).
  \item[Backpropagation] an algorithm that computes the gradient of every parameter by applying the chain rule from the output back to the input (\emph{turunan balik}).
  \item[Backward gradient] see backpropagation (\emph{gradien terbalik}).
  \item[Backward pass differentiable approximation] a stand-in for the gradient on a piece that is not differentiable, using the gradient of the function at the point just after it, so that an attack can still be computed (\emph{BPDA}).
  \item[Barren plateau] a region of the loss landscape of a variational circuit where the gradient shrinks exponentially with the number of qubits (\emph{plateau tandus}).
  \item[Basic feasible solution] a point obtained by setting some variables to zero and then solving for the rest, and still satisfying every constraint (\emph{solusi layak basis}).
  \item[Basis] a set of vectors that is linearly independent and spans the whole space (\emph{basis}).
  \item[Bias-variance tradeoff] the split of model error into error from being too simple and error from being too sensitive to the data (\emph{bias-variansi}).
  \item[Bijection] a perfect pairing of two sets: injective and surjective at once, with nothing left over on either side (\emph{bijeksi}).
  \item[Bootstrap] estimating the distribution of a statistic by resampling from the data itself, many times, with replacement (\emph{bootstrap}).
  \item[Born rule] the rule that turns an amplitude into a probability by squaring its modulus (\emph{aturan Born}).
  \item[Calibration] the agreement between the numbers a model promises and the actual frequency in the group promised the same (\emph{kalibrasi}).
  \item[Catastrophic cancellation] the loss of significant digits when two nearly equal numbers are subtracted, so that only error is left (\emph{pembatalan katastrofik}).
  \item[Cauchy sequence] a sequence whose terms crowd closer and closer together, without needing to name the value they approach (\emph{barisan Cauchy}).
  \item[Central limit theorem] the average of many independent random variables approaches the normal distribution once it is standardized (\emph{Teorema Limit Pusat}).
  \item[Chain rule] the rate of change of a composed function equals the product of the rates at each link in the chain (\emph{aturan rantai}).
  \item[Characteristic polynomial] $\det(A-\lambda I)$; its roots are the eigenvalues of $A$ (\emph{polinomial karakteristik}).
  \item[Coefficient of determination] the share of variation explained by the model, stated as a fraction of the whole; written $R^2$ (\emph{koefisien determinasi}).
  \item[Combination] a choice of several members without regard to their order; there are $\binom{n}{k}$ of them (\emph{kombinasi}).
  \item[Compact] a closed and bounded set; every sequence in it has a subsequence that converges inside that set (\emph{kompak}).
  \item[Completeness axiom] the promise that every nonempty set bounded above in $\mathbb{R}$ has a supremum; this is what leaves the number line with no holes (\emph{aksioma kelengkapan}).
  \item[Condition number] a measure of how much the relative error of the output can grow compared with the relative error of the input (\emph{bilangan kondisi}).
  \item[Conditional probability] the probability of an event given that another event is known to have occurred; written $P(A\mid B)$ (\emph{peluang bersyarat}).
  \item[Confidence interval] a range of values that, if the sampling procedure were repeated, contains the true parameter with a given probability (\emph{interval kepercayaan}).
  \item[Confusion matrix] a four-cell table that compares predictions with reality (\emph{matriks konfusi}).
  \item[Continuous] a function whose graph can be drawn without lifting the pen at that point, that is, $\lim_{x\to c}f(x)=f(c)$ (\emph{kontinu}).
  \item[Convergence in distribution] its distribution function tends to the limit distribution function at every point of continuity; what comes to resemble is the shape of the distribution, not its value (\emph{konvergensi dalam distribusi}).
  \item[Convergence in probability] $P(|X_n-X|>\varepsilon)\to 0$ for every $\varepsilon>0$ (\emph{konvergensi dalam peluang}).
  \item[Convex combination] a weighted mixture of vectors with nonnegative weights that sum to one (\emph{kombinasi konveks}).
  \item[Convex function] a function whose graph always lies below the chord joining two of its points; every local minimum of it is a global minimum (\emph{fungsi konveks}).
  \item[Coordinate vector] the list of coefficients that expresses a vector in a basis (\emph{koordinat}).
  \item[Correlation] a measure of the strength of the linear relation between two variables; it does not state cause and effect (\emph{korelasi}).
  \item[Countable] a set whose members can be numbered with nothing skipped; $\mathbb{Q}$ is countable, $\mathbb{R}$ is not (\emph{terhitung}).
  \item[Countable set] a set whose members can be listed as first, second, third, and so on with nothing left behind (\emph{himpunan terhitung}).
  \item[Counterexample] one example that breaks a statement of the form \"for all\" (\emph{lawan contoh}).
  \item[Covariance] a measure of two quantities moving together, computed from the product of their deviations (\emph{kovarians}).
  \item[Covariance] a measure of how far two variables move together; its unit is the product of their units (\emph{kovariansi}).
  \item[Cross-entropy] a measure of surprise $-\sum_i y_i\log \hat{p}_i$; the more confident the model is on a wrong answer, the larger the penalty (\emph{entropi silang}).
  \item[Cross-validation] assessing a model by dividing the data in turn into training and test parts (\emph{validasi silang}).
  \item[Cumulative distribution function] the function $F(x)=P(X\le x)$; a rising staircase whose height is the probability already accumulated up to $x$ (\emph{fungsi distribusi kumulatif}).
  \item[Data leakage] information from the test data that leaks into the training process without anyone intending it, so that test results look better than reality (\emph{kebocoran data}).
  \item[Data split] the division of data into training, validation, and test sets that do not leak into one another (\emph{pemisahan data}).
  \item[Defective] a matrix whose eigenvalues repeat but whose eigenspace is too narrow to be a basis (\emph{defective}).
  \item[Dense set] a set every point of which can be approached by members of that set other than the point itself (\emph{himpunan padat}).
  \item[Derivative] the limit of the difference quotient; the slope of the tangent line and at the same time the instantaneous rate of change (\emph{turunan}).
  \item[Determinant] one number that measures the factor by which volume changes and at the same time tests invertibility (\emph{determinan}).
  \item[Diagonalization] a change of basis until the matrix becomes diagonal, so that the work is reduced to multiplying numbers (\emph{diagonalisasi}).
  \item[Differential equation] a sentence that ties the rate of change of a quantity to the size of that quantity itself (\emph{persamaan diferensial}).
  \item[Differential privacy] the promise that one person's participation in the data changes the output distribution only slightly, measured by $\varepsilon$ (\emph{privasi diferensial}).
  \item[Digital signature] a value that only the holder of the private key can create but anyone can check with the public key (\emph{tanda tangan digital}).
  \item[Dimension] the number of vectors in a basis; the number belongs to the space, not to one basis that happens to be chosen (\emph{dimensi}).
  \item[Direct sum] the sum of two subspaces that meet only at the zero vector, so that every vector has one way of being written (\emph{jumlah langsung}).
  \item[Double integral] the sum of values over small cells of a region, then its limit as the cells shrink (\emph{integral lipat dua}).
  \item[Duality] formulating an optimization problem from the other side, whose optimal value gives a lower bound for the original problem (\emph{dualitas}).
  \item[Eigenspace] all the eigenvectors for one eigenvalue together with the zero vector; it is a subspace (\emph{ruang eigen}).
  \item[Eigenvalue] a scalar $\lambda$ such that $A\mathbf{v}=\lambda\mathbf{v}$ for some nonzero vector (\emph{nilai eigen}).
  \item[Elementary matrix] an identity matrix that has had one row operation applied to it, so that it becomes a record of that operation (\emph{matriks elementer}).
  \item[Elimination] replacing a system by a simpler system without changing its solution set (\emph{eliminasi}).
  \item[Epsilon-delta definition] the style of defining a limit by naming the output tolerance first, then the input bound that guarantees it (\emph{epsilon-delta}).
  \item[Equilibrium point] a state of a dynamical system that does not change when it is not disturbed (\emph{titik setimbang}).
  \item[Estimator] a rule that turns a sample into a guess at a parameter; the guess it produces is called the estimate (\emph{penaksir}).
  \item[Euler method] a numerical solution of a differential equation that steps along the tangent line; its error is first order in the step (\emph{metode Euler}).
  \item[Euler totient function] the function $\varphi(n)$: the count of numbers between $1$ and $n$ that are coprime with $n$ (\emph{fungsi phi}).
  \item[Expectation] the expected value $\mathbb{E}[X]$: the probability-weighted average, the balance point of a distribution (\emph{harapan}).
  \item[Expectation over transformation] an attack that maximizes the expected loss over the randomness of the defense, usually by averaging gradients over several random copies (\emph{EOT}).
  \item[Exponential distribution] the distribution of waiting time in a process without aging, with density $\lambda e^{-\lambda t}$ (\emph{distribusi eksponensial}).
  \item[Fast gradient sign method] a one-step attack (FGSM) that shifts every coordinate by the budget, in the direction of the sign of the loss derivative (\emph{metode tanda gradien}).
  \item[Field] a set with two operations that is closed, has identities and inverses, and satisfies the distributive law (\emph{lapangan}).
  \item[Floating-point system] binary scientific notation with a fixed number of digits; a number is stored as sign, exponent, and mantissa, so that the spacing between representable numbers is not uniform: dense near zero, sparse at the top end, and two different numbers can fall into the same place (\emph{sistem floating-point}).
  \item[Formal verification] a proof that a program or formula meets its specification, checked by a machine (\emph{verifikasi formal}).
  \item[Functional limit] $f(x)$ can be forced as close as we like to $L$ provided $x$ is close enough to $c$ but not equal to $c$ (\emph{limit fungsi}).
  \item[General solution] the family of all solutions of a differential equation, depending on arbitrary constants (\emph{solusi umum}).
  \item[Generalization] the ability to do problems never seen before, not problems already memorized (\emph{generalisasi}).
  \item[Global sensitivity] the largest shift in the answer to a query that one person can cause in the worst case (\emph{global sensitivity}).
  \item[Gradient] the vector of partial derivatives; the direction of steepest ascent of a function (\emph{gradien}).
  \item[Gradient descent] the parameter update rule that moves against the direction of the gradient (\emph{PGC, penurunan gradien}).
  \item[Harmonic series] the series $1+1/2+1/3+\cdots$; its terms go to zero but its sum diverges (\emph{deret harmonik}).
  \item[Heavy tail] a distribution whose probability of enormous values shrinks too slowly, so that one observation can outweigh all the rest (\emph{ekor berat}).
  \item[Hessian] the matrix of second partial derivatives; the concavity of a function is decided by its sign (\emph{Hessian}).
  \item[Hidden layer] a layer of units that touches neither the input nor the output of the network directly (\emph{lapisan tersembunyi}).
  \item[Hypothesis test] a procedure for deciding, with a stated risk of error, whether the data are enough to reject an initial claim (\emph{uji hipotesis}).
  \item[Implication] a statement of the form if $P$ then $Q$; it is false only when $P$ is true and $Q$ is false (\emph{implikasi}).
  \item[Independent] two events that give each other no news, $P(A\cap B)=P(A)P(B)$; knowing one does not change the probability of the other (\emph{independen}).
  \item[Independent and identically distributed] samples that are mutually independent and share one distribution; the most often used assumption, and the most often violated (\emph{i.i.d.}).
  \item[Induction] the domino pattern of proof: the base case holds, then each case knocks down the next (\emph{induksi}).
  \item[Inner product] a mapping of two vectors to one scalar that measures length and angle; in $\mathbb{R}^n$ its value is $\langle \mathbf{x},\mathbf{y}\rangle=\mathbf{x}^{\top}\mathbf{y}$ (\emph{hasilkali dalam}).
  \item[Integrality gap] the distance between the best value of the LP relaxation and the best value of the integer answer; it measures how far a continuous model misleads (\emph{celah integralitas}).
  \item[Integrating factor] a multiplier that turns the left side of a linear differential equation into a single derivative of a product (\emph{faktor integrasi}).
  \item[Interpolation] finding a function that passes through every data point; different from regression, which does not demand that (\emph{interpolasi}).
  \item[Isomorphism] a bijective mapping that preserves structure; two isomorphic spaces cannot be told apart algebraically (\emph{isomorfisme}).
  \item[Jacobian] the matrix of first partial derivatives of a vector mapping; its determinant measures the local volume scale factor (\emph{Jacobian}).
  \item[Jacobian determinant] the local volume scale factor when coordinates are changed; its absolute value multiplies area or volume in a multiple integral (\emph{determinan Jacobian}).
  \item[Kamus (glossary box)] a small box inside a chapter that explains one new term together with its analogy (\emph{kamus}).
  \item[Karush-Kuhn-Tucker conditions] necessary conditions for optimality in a problem with inequality constraints (\emph{kondisi KKT}).
  \item[Kernel, null space] the set of all vectors mapped to zero by a linear transformation (\emph{inti}).
  \item[Kernel, null space] the set of all inputs mapped to zero by a linear transformation; the same as the null space (\emph{kernel}).
  \item[Lagrange multiplier] an auxiliary variable that puts a price on a constraint; its value shows how much the objective could improve if the constraint were loosened (\emph{pengali Lagrange}).
  \item[Law of large numbers] the promise that the average of many observations sticks to the true mean, with the chance of missing it shrinking (\emph{hukum bilangan besar}).
  \item[Least squares] the choice of parameters that minimizes the sum of squared residuals when there is no exact solution (\emph{kuadrat terkecil}).
  \item[Linear map] a mapping that preserves addition and scalar multiplication (\emph{transformasi linear}).
  \item[Linear program] the problem of maximizing or minimizing a linear function with linear constraints (\emph{program linear}).
  \item[Lipschitz condition] the condition that the rate of change of a function is bounded by one constant; this is what guarantees a unique solution of a differential equation (\emph{Lipschitz}).
  \item[Logical connective] an operation that combines statements: and, or, not, if-then, if-and-only-if (\emph{penghubung logika}).
  \item[Logistic regression] a model for binary output that maps a linear combination to a probability through the logistic function (\emph{regresi logistik}).
  \item[Loss function] a number we want to make small while the model is trained; it measures how wrong a prediction is, not how right (\emph{fungsi kerugian}).
  \item[LP relaxation] replacing the integer condition by a real condition so that an integer problem can be solved as an LP (\emph{relaksasi LP}).
  \item[Lyapunov function] a positive-valued function that decreases along the trajectories of a system; its existence proves stability without solving the equations (\emph{fungsi Lyapunov}).
  \item[Machine epsilon] the distance between $1$ and the next machine number; the size of one step at scale one (\emph{epsilon mesin}).
  \item[Modular inverse] a number $b$ such that $ab\equiv 1 \pmod n$; the undo button of clock arithmetic (\emph{invers modular}).
  \item[Monotone sequence] a sequence that only rises or only falls, never turning back (\emph{barisan monoton}).
  \item[Multicollinearity] the situation in which explanatory variables are nearly linearly dependent on one another, so that regression coefficients become unstable (\emph{multikolinearitas}).
  \item[Multiplicity] the number of times an eigenvalue occurs as a root of the characteristic polynomial, which is not always equal to the dimension of its eigenspace (\emph{multiplisitas}).
  \item[Necessary condition] a condition that must hold for a statement to be true; without it the statement is impossible, but meeting it is not yet enough (\emph{syarat perlu}).
  \item[Nested cross-validation] two layers of folds: the outer layer assesses, the inner layer selects, so that model selection does not leak test data (\emph{validasi silang bersarang}).
  \item[Network intrusion detection system] a system that watches network traffic and decides whether a flow of packets is an attack (\emph{NIDS}).
  \item[Norm] a measure of the length of a vector; the $L_\infty$ norm takes the largest component (\emph{norma}).
  \item[Normal distribution] the symmetric bell distribution determined entirely by its mean and its variance (\emph{distribusi normal}).
  \item[Null hypothesis] the initial claim under test, usually \"there is no effect\"; it is rejected only when the data fit it very badly (\emph{hipotesis nol}).
  \item[Obfuscated gradient] a gradient handed to the attacker that does not point in the right direction, usually because a piece is not differentiable or is random (\emph{gradien dikaburkan}).
  \item[Odds ratio] the ratio of the probability of success to the probability of failure; the exponential of a logistic regression coefficient states how many times that ratio multiplies per unit (\emph{rasio peluang}).
  \item[Ordered field] a field whose members can be ordered and whose arithmetic rules respect that order (\emph{lapangan terurut}).
  \item[Ordinary differential equation] an equation that relates a function to its derivative with respect to one independent variable (\emph{persamaan diferensial biasa}).
  \item[Orientation] an agreement about which direction of rotation counts as positive; reversing it reverses the sign of an integral (\emph{orientasi}).
  \item[Orthogonal] mutually perpendicular in the sense that their inner product is zero (\emph{ortogonal}).
  \item[Orthogonal matrix] a square matrix whose columns are orthonormal; it preserves length and angle (\emph{matriks ortogonal}).
  \item[Orthonormal] an orthogonal set each of whose members has length one (\emph{ortonormal}).
  \item[Out-of-sample] data not used when the model was trained, so that they are honest for judging it (\emph{luar sampel}).
  \item[Overfitting] memorizing the training data including its noise, so that performance on new data falls (\emph{overfitting}).
  \item[Padding] extra content added to a packet before it is sent so that every packet has a fixed size, so that packet size no longer leaks its content (\emph{bantalan}).
  \item[Partition] a division of a set into pieces that do not overlap and that cover the whole set (\emph{partisi}).
  \item[Perceptron] the simplest linear classifier: the sign of the inner product of the weights and the input (\emph{perceptron}).
  \item[Poles and zeros] values of a complex variable at which a function blows up or vanishes; used in the analysis of dynamical systems and control (\emph{kutub dan nol}).
  \item[Polynomial time] the class of algorithms whose number of steps grows no faster than a polynomial in the size of the input (\emph{waktu polinomial}).
  \item[Power series] a series of the form $\sum a_n(x-c)^n$; it holds only inside its radius of convergence (\emph{deret pangkat}).
  \item[Power set] the set of all subsets, including the empty set and the set itself; its cardinality is $2^n$ (\emph{himpunan kuasa}).
  \item[Prediction interval] an interval for one observation that has not yet occurred; always wider than the interval for a mean (\emph{selang prediksi}).
  \item[Primitive root] a member that, raised repeatedly to powers, passes through every other member of the cyclic group before returning to one; in cryptography, a starting position that reaches every position of the key (\emph{akar primitif}).
  \item[Privacy budget] the amount $\varepsilon$ a system may spend; every query or training step draws down what is left (\emph{anggaran privasi}).
  \item[Probability] a weight on every possibility, nonnegative and summing to one (\emph{peluang}).
  \item[Probability amplitude] a complex number whose squared modulus gives a probability; different from a probability because two amplitudes can cancel each other (\emph{amplitudo}).
  \item[Problem space] the space of real objects an attacker may change, for example packets and timing; different from the feature space in which the model works (\emph{ruang masalah}).
  \item[Projection] the point in a subspace closest to a vector, that is, its perpendicular image (\emph{proyeksi}).
  \item[Proof term] an expression that is at once a proof and a value; checking a proof is the same as checking a type (\emph{term bukti}).
  \item[Proposition] a sentence that is either true or false, not a question or a command; in this book also the name of a theorem environment (\emph{proposisi}).
  \item[Pseudoinverse] a generalized inverse $A^{+}$ that still works when some directions have been flattened: what can be inverted is inverted, what is lost is left alone (\emph{pseudoinvers}).
  \item[P-value] the probability of obtaining a result at least as extreme as the one observed if the null hypothesis is true; not the probability that the null hypothesis is true (\emph{nilai-p}).
  \item[Quantifier] a marker for all ($\forall$) or there exists ($\exists$); the order in which they are used changes the meaning of a statement (\emph{kuantor}).
  \item[Quantization] replacing every value by the nearest multiple of a fixed step; the result is a staircase function whose derivative is zero almost everywhere (\emph{kuantisasi}).
  \item[Quantum computing era] the era when quantum computers are large enough to run Shor's algorithm on the keys used today (\emph{zaman kuantum}).
  \item[Radius of convergence] the greatest distance from the center at which a power series still converges (\emph{jari-jari konvergensi}).
  \item[Random variable] a function that maps the outcome of an experiment to a number (\emph{variabel acak}).
  \item[Rank] the number of independent directions that survive after mapping, that is, $\dim\operatorname{im}T$ (\emph{rank}).
  \item[Rank-nullity theorem] $\dim\ker T+\dim\operatorname{im}T=\dim V$ for a linear transformation on a finite-dimensional space (\emph{rank-nullity}).
  \item[Ratio test] a series converges absolutely if $\lim|a_{n+1}/a_n|<1$, and diverges if that limit is greater than one (\emph{uji rasio}).
  \item[Rearrangement] swapping the order of the terms of a series; in a series that converges absolutely the sum does not change, in one that does not, the sum can change (\emph{susunan ulang}).
  \item[Regression] drawing the function closest to all the points, not the function that passes through all of them (\emph{regresi}).
  \item[Regression line] the line (or plane) that minimizes the sum of squared residuals; equivalent to the orthogonal projection onto the column space (\emph{garis regresi}).
  \item[Regularization] a penalty on the size of the parameters so that the model is not too sensitive to the training data (\emph{regularisasi}).
  \item[Residual] the difference between an observed value and the value predicted by the model (\emph{residu}).
  \item[Residue class] the set of numbers that leave the same remainder when divided by $n$, like three in the morning and three in the afternoon (\emph{kelas sisa}).
  \item[Riemann integrability] the property of a function whose upper and lower sums can be squeezed until their difference vanishes (\emph{keterintegralan}).
  \item[Riemann integral] the limit of the sum of rectangle areas as the width of the partition shrinks (\emph{integral Riemann}).
  \item[ROC curve] a plot of the true rate against the false positive rate over all thresholds; the area under it is called AUROC (\emph{kurva ROC}).
  \item[Rounding] replacing a real number by the nearest machine number (\emph{pembulatan}).
  \item[Rounding error] the difference between a real number and its machine representative, which piles up in long sums (\emph{kesalahan pembulatan}).
  \item[Rule of three] if no event occurs in $n$ trials, the event rate may still be as high as about $3/n$ at 95 percent confidence (\emph{aturan tiga}).
  \item[Runge-Kutta] a family of numerical methods that evaluate the derivative several times per step for a smaller error (\emph{Runge-Kutta}).
  \item[Sample complexity] the number of examples needed for a model to generalize with a given confidence (\emph{kompleksitas sampel}).
  \item[Sensitivity] in differential privacy, the largest change in output that one person can cause (\emph{sensitivitas}).
  \item[Sequence] a row of numbers given a running order $a_1,a_2,\dots$; the order is part of the definition, not something added later (\emph{barisan}).
  \item[Series] the sum of the terms of a sequence without stopping, whose value is read as the limit of the partial sums (\emph{deret}).
  \item[Similar] two matrices that represent the same linear mapping in two different bases, $B=P^{-1}AP$ (\emph{serupa}).
  \item[Singular value] the square root of an eigenvalue of $A^{\top}A$; it measures the length of a principal axis of the mapping by the matrix $A$ (\emph{nilai singular}).
  \item[Singular value decomposition] writing a matrix as $U\Sigma V^{\top}$ with $U,V$ orthogonal (\emph{dekomposisi nilai singular, SVD}).
  \item[Slack variable] an extra variable that records the unused remainder, so that a no more than constraint becomes an equation (\emph{variabel kendur}).
  \item[Softmax] a mapping of a row of scores into probabilities that sum to one, $e^{z_i}/\sum_j e^{z_j}$ (\emph{softmax}).
  \item[Sorry] a Lean~4 marker that a proof has been postponed; the machine accepts it but records it as an axiom (\emph{sorry}).
  \item[Stability] the property of an equilibrium point or dynamical system that stays close to its state after a small disturbance (\emph{kestabilan}).
  \item[Standard error] the standard deviation of an estimator, not the standard deviation of the data; it measures the wobble of one guess (\emph{galat baku}).
  \item[Standard form] a uniform way of writing an optimization problem: one objective, constraints in the form no more than zero, and nonnegative variables (\emph{bentuk baku}).
  \item[Stochastic gradient descent] gradient descent that uses a small fraction of the data at every step; cheaper and noisier (\emph{penurunan gradien stokastik}).
  \item[Strong induction] induction that may use every earlier case, not only the nearest one (\emph{induksi kuat}).
  \item[Subsequence] a sequence obtained by discarding some terms without changing the order of the remaining terms (\emph{subbarisan}).
  \item[Subspace] a subset of a vector space that is closed under addition and scalar multiplication (\emph{subruang}).
  \item[Summary box] a box of bullet points at the end of a chapter that can be read on its own without the body of the chapter (\emph{ringkasan}).
  \item[Supremum, least upper bound] the least upper bound; the smallest number that is not smaller than every member of the set (\emph{supremum}).
  \item[Tactic] a step-by-step command in Lean~4 that changes the goal of a proof (\emph{taktik}).
  \item[Taylor series] writing a function as a sum of derivative terms at one point, together with a remainder that measures the truncation error (\emph{deret Taylor}).
  \item[Threat model] the complete list of an attacker's abilities: the size of the disturbance, what the attacker knows, and what the attacker wants (\emph{model ancaman}).
  \item[Transferability] the property of an adversarial example made for one model that also fools another model trained separately (\emph{transferabilitas}).
  \item[Truncation error] the difference that appears because an infinite process is replaced by a finite piece (\emph{galat pemotongan}).
  \item[Type class] a list of abilities a type must have in order to be usable in a given place (\emph{typeclass}).
  \item[Uniform continuity] continuity with one $\delta$ that works for every point of the set, not one $\delta$ per point (\emph{kekontinuan seragam}).
  \item[Uniform convergence] a sequence of functions that approaches its limit function at the same rate over the whole domain (\emph{kekonvergenan seragam}).
  \item[Uniform convergence] one rate for the whole domain, not one schedule per point (\emph{konvergensi seragam}).
  \item[Uniqueness] the guarantee that one initial state gives rise to only one solution trajectory (\emph{ketunggalan}).
  \item[Unit in the last place] the smallest step a machine can distinguish at the scale of one number (\emph{ULP}).
  \item[Unitary] a matrix that preserves the inner product, $U^{*}U=I$; the complex version of an orthogonal matrix (\emph{uniter}).
  \item[Upper limit] the largest value that a sequence still visits arbitrarily often near its end, written $\limsup_{n\to\infty}a_n$ (\emph{limit superior}).
  \item[Upper sum] a deliberately enlarged estimate of an integral, using the largest value on each piece (\emph{jumlah atas}).
  \item[Variance] a measure of spread around the mean; its square root is the standard deviation (\emph{variansi}).
  \item[Variance inflation factor] how much the variance of a coefficient is inflated by other explanatory variables carrying similar information; abbreviated VIF (\emph{varians inflasi}).
  \item[Vector space] a set with addition and scalar multiplication that is closed and satisfies eight axioms (\emph{ruang vektor}).
  \item[Vertex] a corner of the feasible region, that is, a point that is never the midpoint between two other feasible points (\emph{verteks}).
  \item[Wilson interval] a confidence interval for a proportion found by inverting a test, so that it stays sensible when the proportion approaches zero or one (\emph{interval Wilson}).
  \item[Wilson interval] a confidence interval for a proportion that stays sensible when the proportion approaches zero or one (\emph{Wilson}).
\end{description}
